\chapter{Pétrole brut}\label{ch:m3:crude-oil}

À 14~h~30, heure de New York, le lundi 20 avril 2020, le contrat de mai sur le West Texas
Intermediate a réglé à moins 37.63 dollars le baril. La séance avait
ouvert le dimanche soir à plus 17.73. Le contrat n'avait jamais été négocié sous zéro
depuis sa cotation en 1983, et il a expiré le lendemain à plus 10.01. Pendant environ
vingt minutes, cet après-midi-là, le détenteur du contrat devait payer plusieurs dizaines de dollars
le baril pour se débarrasser d'un pétrole qu'il aurait autrement dû réceptionner, dans une ville de
l'Oklahoma dont les réservoirs étaient remplis aux trois quarts et dont l'espace restant était déjà
réservé. Ce chapitre explique ce qu'est un \omterm{def:m3:crude-oil:benchmark}{brut de référence}, comment sont construits les trois qui
fixent l'essentiel du prix du pétrole mondial, et pourquoi le stockage qui se trouve derrière un
indice de référence compte autant que le pétrole qu'il contient.

\section{Qualité des bruts et indices de référence}

Le pétrole brut n'est pas un produit unique mais des centaines de qualités, chacune étant un mélange
d'hydrocarbures dont la valeur pour une raffinerie dépend de ce qu'elle peut en tirer
(\cref{ch:m3:refined-products-and-cracks}). Deux grandeurs résument l'essentiel de cette
valeur.

\begin{definition}[Densité API]\label{def:m3:crude-oil:api}
La \emph{densité API}\index{densité API} d'un pétrole brut est
\[
\text{°API} = \frac{141.5}{\mathrm{SG}} - 131.5 ,
\]
où $\mathrm{SG}$ est sa densité à \qty{60}{\degree F} par rapport à l'eau.
L'eau a 10 degrés API~; plus le pétrole est léger (moins dense), plus sa
densité API est élevée.
\end{definition}

\begin{definition}[Brut doux, brut soufré]\label{def:m3:crude-oil:sweet}
Un pétrole brut est un \emph{brut doux}\index{brut doux} si sa teneur en soufre est
faible, par convention inférieure à 0.5~\% en masse, et un \emph{brut soufré}\index{brut soufré}
si elle est supérieure. (Les statistiques de l'Energy Information Administration américaine
retiennent un seuil de 1~\%.)
\end{definition}

Les \omterm{def:m3:crude-oil:sweet}{bruts doux} et légers donnent davantage d'essence et de gazole avec moins de traitement et
valent davantage~; les bruts lourds et soufrés exigent des raffineries plus complexes. Un brut de
densité 0.827 a une \omterm{def:m3:crude-oil:api}{densité API} de 39.6 degrés. Les contrats physiques corrigent la qualité au moyen
d'échelles d'ajustement~: tant par degré API au-dessus d'une référence, tant en moins par dixième
de pour cent de soufre. Avec \qty{0.04}{\$/\barrel} par degré et
\qty{0.12}{\$/\barrel} par 0.1~\% de soufre, un brut de 42° à 0.25~\% de soufre vaut
\qty{0.34}{\$/\barrel} de plus qu'une référence de 38° à 0.40~\%.

\begin{definition}[Brut de référence]\label{def:m3:crude-oil:benchmark}
Un \emph{brut de référence}\index{brut de référence} est un pétrole brut (ou un panier de
bruts) dont le prix est évalué ou négocié avec assez de volume et de transparence pour que
les prix des autres qualités soient fixés comme des écarts par rapport à lui.
\end{definition}

Trois familles d'indices de référence fixent l'essentiel du prix du pétrole mondial~: le complexe
de la mer du Nord, construit autour du \omterm{def:m3:crude-oil:brent}{Dated Brent}, le baril américain de l'intérieur des terres, le WTI livré à Cushing,
dans l'Oklahoma, et les références du Moyen-Orient, Dubaï et Oman, auxquelles s'ajoute désormais le Murban.

\section{Le complexe de la mer du Nord}

L'indice de référence maritime n'est pas un prix unique mais un ensemble de marchés liés, depuis les
cargaisons physiques chargées dans les semaines à venir jusqu'aux contrats à terme portant sur des années
(\cref{fig:m3:crude-oil:complex}).

\begin{definition}[Dated Brent, Forward Brent, Brent CFD]\label{def:m3:crude-oil:brent}
Le \emph{Dated Brent}\index{Dated Brent} est l'évaluation quotidienne de l'\omterm{def:m3:physical-commodity-markets:pra}{agence de cotation des prix} (PRA)
des cargaisons physiques de la mer du Nord dont les dates de chargement sont déjà fixées dans les
semaines à venir, établie au niveau de la qualité la plus concurrentielle d'un panier de bruts, après
ajustements de qualité. Le \emph{Forward Brent}\index{Forward Brent} est le marché des contrats à terme
portant sur une cargaison de l'une des qualités du panier chargée à un mois futur déterminé,
dont la qualité et les dates de chargement sont désignées ultérieurement par le vendeur. Un \emph{Brent
CFD}\index{Brent CFD} est un swap à court terme qui échange la moyenne du Dated Brent
sur une semaine contre un prix fixe convenu par rapport au Forward Brent.
\end{definition}

\begin{omfigure}
\centering
\begin{tikzpicture}[font=\small,
  lay/.style={draw,thick,rounded corners=2pt,align=center,minimum height=0.9cm,text width=3.7cm},
  arr/.style={{Stealth[length=2.2mm]}-{Stealth[length=2.2mm]},thick}]
\node[lay,draw=omDef,fill=omDef!8] (dated) at (0,0) {Dated Brent\\cargaisons physiques,\\dates fixées};
\node[lay,draw=omProp,fill=omProp!8] (cfd) at (4.9,0) {Brent CFD\\swaps hebdomadaires};
\node[lay,draw=omThm,fill=omThm!8] (fwd) at (9.8,0) {Forward Brent\\cargaison à un mois donné};
\node[lay,draw=omExo,fill=omExo!8] (fut) at (9.8,-2.2) {future Brent de l'ICE\\réglé en espèces sur un indice};
\draw[arr] (dated) -- (cfd);
\draw[arr] (cfd) -- (fwd);
\draw[arr] (fwd) -- node[right,font=\footnotesize,align=left] {exchange for\\physical (EFP)} (fut);
\node[font=\footnotesize,align=center] at (3.2,-2.2) {délai avant le chargement croît de gauche à droite~:\\jours et semaines \quad$\longrightarrow$\quad mois et années};
\end{tikzpicture}

\omcaption{Le complexe Brent. Les CFD fixent l'écart entre la cargaison physique
des semaines à venir et le mois à terme~; un exchange for physical transforme une position
sur contrats à terme en cargaison à terme. Schéma simplifié.}\label{fig:m3:crude-oil:complex}
\end{omfigure}

Le panier existe parce que la production du gisement de Brent d'origine a décliné il y a longtemps~;
plus il y a de qualités, plus l'indice de référence est difficile à manipuler. Chaque qualité entre
avec une prime ou une décote de qualité, et chaque jour le prix net le plus bas fixe
l'évaluation. Si un jour le Forties est offert à \qty{81.05}{\$/\barrel} et que les autres
qualités, nettes de leurs primes, le sont à \qty{81.20}{\$/\barrel} ou plus, le \omterm{def:m3:crude-oil:brent}{Dated Brent}
vaut \qty{81.05}{\$/\barrel}.

\begin{dated}{2026-09}{Le panier}\label{dat:m3:crude-oil:basket}
Le \omterm{def:m3:crude-oil:brent}{Dated Brent} de Platts reflète le Brent, le Forties, l'Oseberg, l'Ekofisk et le Troll et, depuis les
livraisons de juin 2023, le WTI Midland, qualité américaine livrée à Rotterdam, avec des évaluations
établies sur une base FOB, pour des cargaisons de 700\,000 barils (contre 600\,000 auparavant)
chargées de dix jours à un mois à l'avance. Platts décrit cet indice comme le point de référence
pour environ les deux tiers du commerce mondial de pétrole. Les contrats à terme Brent de l'ICE sont
réglés en espèces sur la base de l'indice Brent de l'ICE, avec un passage à la livraison physique vers des cargaisons à terme.
\end{dated}

\section{L'indice de référence américain de l'intérieur des terres}

Le WTI est un brut acheminé par pipeline, et son contrat à terme est livré physiquement en un
seul point de livraison intérieur~: Cushing, dans l'Oklahoma, une ville de parcs de réservoirs où les pipelines issus des
bassins producteurs rejoignent les canalisations vers les raffineries. Comme les contrats à terme y sont livrés,
le stockage à Cushing fait partie du contrat.

\begin{definition}[Dernier jour de négociation]\label{def:m3:crude-oil:ltd}
Le \emph{dernier jour de négociation}\index{dernier jour de négociation} d'un contrat à terme est le
dernier jour où il peut être négocié. Une position encore ouverte à sa clôture doit être
réglée par livraison (ou, pour un contrat réglé en espèces, au prix de
règlement-livraison final).
\end{definition}

\begin{dated}{2026-09}{Les contrats à terme sur le WTI}\label{dat:m3:crude-oil:cl}
Contrats à terme NYMEX Light Sweet Crude Oil~: 1\,000 barils, livraison \omterm{def:m3:physical-commodity-markets:incoterms}{franco à bord} (FOB) sur tout
point de raccordement à un pipeline ou toute installation de stockage à Cushing, dans l'Oklahoma. La négociation prend fin trois
jours ouvrés avant le 25e jour calendaire du mois précédant le mois de livraison (trois
jours ouvrés avant le dernier jour ouvré précédant le 25e si celui-ci n'est pas un jour
ouvré). Le mois actif passe au contrat suivant quelques jours avant l'échéance, et les
autres mois se règlent ensuite par des transactions en calendar spread.
\end{dated}

La plupart des bruts américains acheminés par pipeline ne se fixent pas non plus sur un seul jour.

\begin{definition}[Moyenne du mois calendaire]\label{def:m3:crude-oil:cma}
Un prix de \emph{moyenne du mois calendaire}\index{moyenne du mois calendaire} (CMA) est la
moyenne des règlements-livraisons quotidiens ou des évaluations d'un indice de référence sur les
jours ouvrés d'un mois calendaire.
\end{definition}

Un baril livré dans un mois et dont le prix est fixé sur la CMA du contrat à terme est exposé, chaque
jour de ce mois, à l'échéance la plus proche, qui change d'identité lorsqu'il
arrive à échéance au milieu du mois~: les formules de CMA sont donc souvent accompagnées
d'un ajustement «~roll~» pour la part du mois fixée sur le second contrat.

\section{Tarification au Moyen-Orient}

Les grands producteurs du Golfe vendent l'essentiel de leur brut par contrats à long terme aux raffineurs,
et non sur un marché au comptant, à des prix qu'ils fixent eux-mêmes.

\begin{definition}[Prix de vente officiel]\label{def:m3:crude-oil:osp}
Un \emph{prix de vente officiel}\index{prix de vente officiel} (OSP) est le
différentiel mensuel publié par une compagnie pétrolière nationale pour chaque qualité et chaque
zone de destination, à ajouter à une moyenne de référence désignée sur le mois de chargement pour
fixer le prix de ses cargaisons à terme.
\end{definition}

Pour l'Asie, la référence est depuis longtemps la moyenne des évaluations de Dubai et d'Oman~;
pour août 2025, Saudi Aramco a fixé son Arab Light pour l'Asie à \qty{2.20}{\$/\barrel}
au-dessus de cette moyenne. En mars 2021, une bourse de contrats à terme d'Abou Dhabi a commencé
à négocier les premiers contrats à terme sur le Murban, un brut léger (\omterm{def:m3:crude-oil:api}{densité API} de 40°) livré physiquement à
Fujairah, et en juin 2021 le producteur d'Abou Dhabi a rattaché son OSP du Murban à ce
contrat.

\begin{dated}{2026-09}{Le Golfe en 2026}\label{dat:m3:crude-oil:hormuz}
Après les frappes aériennes américaines et israéliennes sur l'Iran le 28 février 2026, les flux de brut et
de produits transitant par le détroit d'Ormuz ont chuté, selon les mots de l'Agence internationale
de l'énergie (AIE), «~d'environ 20~mb/j avant la guerre à un filet~». L'AIE l'a qualifiée de la
plus grande perturbation de l'offre de l'histoire du marché pétrolier mondial. Les contrats à terme Brent
se sont négociés «~à un cheveu de 120~\$/bbl~» en mars~; les membres de l'AIE ont convenu le 11 mars de
mettre à disposition 400 millions de barils de stocks d'urgence. Le détroit a été
«~de nouveau fermé dans les faits au début de juillet~»~; le \omterm{def:m3:crude-oil:brent}{Dated Brent} a terminé juillet à
\qty{96.80}{\$/\barrel}. L'indice de référence maritime a grimpé bien au-dessus de l'indice intérieur
(la \cref{fig:m3:physical-commodity-markets:differential}).
\end{dated}

\section{Le 20 avril 2020 pas à pas}

Au premier trimestre 2020, un marché déjà en surproduction a perdu une grande part de
sa demande à cause de la pandémie. Les stocks américains de brut ont augmenté de 15.8~millions de barils par semaine pendant
quatre semaines. Cushing, dont la capacité utile de stockage est de 75.8~millions de barils, est passé
d'un taux de remplissage d'environ la moitié à la mi-mars à environ 76~\% le 17 avril, et la presse
a rapporté que l'essentiel de l'espace restant était engagé. Un détenteur du contrat de mai à l'échéance
devait prendre livraison de 1\,000 barils par contrat à Cushing, ce qui exigeait de la place dans les
réservoirs ou de la capacité de pipeline réservée à l'avance~; la plupart des détenteurs qui n'étaient
pas dans ce métier ont dû vendre avant la clôture du 21 avril.

\begin{proposition}[Plancher de stockage]\label{prop:m3:crude-oil:floor}
Si un trader peut prendre livraison du contrat le plus proche, stocker le pétrole pendant $m$ mois au
coût $c$ par baril-mois et le livrer dans le contrat éloigné, alors, en négligeant
le financement, le prix du contrat proche $F_1$ et celui du contrat éloigné $F_2$ vérifient
$F_1 \ge F_2 - c\,m$. Le spread $F_2 - F_1$ ne peut dépasser le coût du stockage que
si le stockage n'est pas disponible à ce coût.
\end{proposition}

\begin{proof}
Sinon, acheter le contrat proche, en prendre livraison, le stocker et vendre le contrat éloigné~:
les flux de trésorerie valent $-F_1 - cm + F_2 > 0$, sans aucun risque.
\end{proof}

Le 20 avril, le contrat de juin a clôturé à \qty{20.43}{\$/\barrel}~: le spread mai--juin
est passé de $-\$6.76$ le vendredi à $-\$58.06$, le plus large de l'historique quotidien de l'EIA
depuis 1985 (le deuxième plus large est de 8.49~\$). Un trader qui aurait pu stocker un
baril à Cushing pendant un mois aurait touché 58.06~\$ pour cela, alors que le stockage que CME Group
signale avoir coûté environ 50~cents le baril lors d'une enchère de stockage à Cushing en mars. Personne disposant
d'un accès à des réservoirs vides n'est resté pour prendre la transaction, et le plancher de
la \cref{prop:m3:crude-oil:floor} n'a pas joué.

\begin{omfigure}
\begin{tikzpicture}
\begin{axis}[width=11cm,height=5cm,
  xlabel={jours calendaires depuis le 2 mars 2020},ylabel={\$/bbl},
  tick label style={font=\small},label style={font=\small},
  xmin=0,xmax=90,ymin=-45,ymax=55,grid=major,grid style={gray!25},
  legend columns=2,legend style={font=\footnotesize,at={(0.5,-0.3)},anchor=north,draw=none,fill=none,/tikz/every even column/.append style={column sep=8pt}},
  table/col sep=comma]
\addplot[omDef,very thick,mark=*,mark size=1pt] table[x=day,y=c1]{figdata/markets-3/02-crude-oil/spring2020.csv};
\addplot[omThm,very thick,dashed] table[x=day,y=c2]{figdata/markets-3/02-crude-oil/spring2020.csv};
\draw[gray] (axis cs:0,0) -- (axis cs:90,0);
\node[font=\footnotesize,anchor=west] at (axis cs:51,-37.63) {$-37.63$ le 20 avril};
\legend{échéance la plus proche (mai jusqu'au 21 avril),échéance suivante}
\end{axis}
\end{tikzpicture}

\omcaption{Contrats à terme WTI, règlements quotidiens, de mars à mai 2020. Le contrat le plus proche
est passé sous le second à mesure que les stocks se remplissaient~; la veille de son expiration, il a réglé à
$-\$37.63$ tandis que le contrat de juin réglait à 20.43~\$. Données~: US Energy Information
Administration, contrats 1 et 2 du NYMEX.}\label{fig:m3:crude-oil:spring}
\end{omfigure}

\begin{omfigure}
\centering
\begin{tikzpicture}[font=\footnotesize]
\draw[thick,-{Stealth[length=2.5mm]}] (0,0) -- (13,0);
\foreach \x/\t/\l in {0.3/{dim. 18:00}/{ouverture de la séance, 17.73~\$},3.0/{lun. 13:00}/{début de la chute},5.2/{lun. 14:08}/{0~\$},7.4/{lun. 14:29}/{plus bas, $-40.32$~\$},9.4/{lun. 14:30}/{règlement à $-37.63$~\$},12.0/{mar. 14:30}/{expiration, 10.01~\$}}{
  \draw[thick] (\x,0.12) -- (\x,-0.12);
  \node[anchor=south,align=center] at (\x,0.2) {\t};
  \node[anchor=north,align=center,text width=2cm] at (\x,-0.2) {\l};}
\end{tikzpicture}

\omcaption{Les deux dernières séances du contrat WTI de mai 2020, heure de New York. Le
règlement à 14:30 est le prix moyen des transactions de spread calendaire réalisées dans les deux
minutes précédentes, appliqué au contrat de juin. Le prix est passé de 0~\$ à $-\$39.55$
entre 14:08 et 14:28. Source~: CFTC, rapport intermédiaire du personnel,
novembre 2020.}\label{fig:m3:crude-oil:timeline}
\end{omfigure}

Trois caractéristiques du contrat ont transformé un excédent en prix négatif. Il était
livrable physiquement en un seul lieu dont les capacités de stockage étaient presque saturées. Les positions
ouvertes sur le contrat de mai, 634\,727 contrats au début d'avril, étaient encore de
108\,593 contrats au début de la séance du 20 avril, un niveau plus élevé que d'ordinaire
pour l'avant-dernier jour. Enfin, le règlement, établi à partir de transactions de spread sur une
fenêtre de deux minutes, est intervenu à la fin de la séance au cours de laquelle ces positions étaient clôturées.
Les fonds détenant des contrats de l'échéance la plus proche s'en sont détournés~: le United States Oil Fund,
un fonds indiciel coté sur pétrole qui détenait habituellement le contrat le plus proche, a réparti
ses avoirs sur des échéances plus éloignées dans les jours entourant l'événement. La bourse avait averti
ses membres, plus tôt en avril, que les prix pouvaient devenir négatifs, et elle a remplacé
ses modèles d'options par le modèle de Bachelier (livre~2, chapitre~13), qui autorise
des prix sous-jacents négatifs. \Cref{fig:m3:crude-oil:spread} montre à quel point la journée
était exceptionnelle~: sur 39 ans de données quotidiennes, le spread a dépassé 5~\$ seulement
25 jours.

\begin{omfigure}
\begin{tikzpicture}
\begin{axis}[width=11cm,height=4.8cm,
  ylabel={second moins premier, \$/bbl},
  tick label style={font=\small},label style={font=\small},
  xmin=1985,xmax=2024.5,ymin=-8,ymax=11,grid=major,grid style={gray!25},
  xtick={1985,1990,1995,2000,2005,2010,2015,2020},xticklabel style={/pgf/number format/1000 sep={}},
  table/col sep=comma]
\addplot[omDef,thin] table[x=t,y=spread]{figdata/markets-3/02-crude-oil/spread.csv};
\draw[gray] (axis cs:1985,0) -- (axis cs:2024.5,0);
\node[font=\footnotesize,anchor=west] at (axis cs:1986,8.3) {20 avril 2020 : 58.06~\$ (hors échelle)};
\end{axis}
\end{tikzpicture}

\omcaption{Second contrat WTI moins contrat le plus proche, de 1985 à avril 2024~: les mercredis,
plus chaque jour où l'écart dépasse 10~\$, plafonné à 10~\$. Les valeurs positives correspondent au contango. Données~: US
Energy Information Administration, contrats 1 et 2 du NYMEX.}\label{fig:m3:crude-oil:spread}
\end{omfigure}

\section{Tutoriel~: le calendrier d'expiration et le plancher de stockage}\label{tut:m3:crude-oil:floor}

\begin{tutorial}
\textbf{Objectif.} Calculer le \omterm{def:m3:crude-oil:ltd}{dernier jour de négociation} de n'importe quel contrat WTI et le coût de
stockage que le marché impliquait le 20 avril 2020. \textbf{État final~:}
\cref{fig:m3:crude-oil:spring,fig:m3:crude-oil:spread}, expiration le 21 avril 2020 et
un coût de stockage implicite de 58.06~\$ par baril-mois.
\begin{tutsteps}
\item \textbf{Le calendrier.} Jours ouvrés, puis la règle d'expiration de l'encadré « Dated ».
\omcode{code/firm/crude/firm_crude.py}{33}{54}{Jours ouvrés et dernier jour de négociation du WTI.}\label{lst:m3:crude-oil:ltd}
\item \textbf{Le plancher.} \Cref{prop:m3:crude-oil:floor} et son inverse.
\omcode{code/firm/crude/firm_crude.py}{65}{74}{Le plancher de stockage et le coût de stockage implicite.}\label{lst:m3:crude-oil:floor}
\item \textbf{Exécuter} \texttt{m3\_crude.april\_20()}, \texttt{spread\_stats()} et
\texttt{fig\_crude.py} sur la série EIA des contrats 1 à 4.
\end{tutsteps}
\textbf{Pour aller plus loin.} Ajouter le calendrier des jours fériés de la bourse et lister chaque
expiration depuis 2015~; calculer le coût de stockage implicite à partir des contrats 2 et 3 sur les mêmes
semaines et voir à quel point il était moins extrême.
\end{tutorial}

\section{Construction~: la boîte à outils du brut}\label{bld:m3:crude-oil:crude}

\begin{build}
\bfield{Rôle} La mini-entreprise valorise le brut physique par rapport à trois familles de
références~: elle a besoin de l'arithmétique de la qualité, de la règle du grade le plus compétitif, du
calendrier d'expiration des contrats à terme qu'elle couvre, des \omterm{def:m3:crude-oil:cma}{moyennes du mois calendaire} et du plancher
de stockage qui indique quand un spread rémunère des cuves.
\bfield{Interface} \texttt{api\_gravity(sg)},
\texttt{specific\_gravity(api)}~; \texttt{quality\_adjusted(price, api, sulphur, ref\_api, ref\_sulphur, per\_api,
per\_tenth\_sulphur)}~; \texttt{basket\_benchmark(grade\_prices, premiums)}~;
\texttt{cl\_last\_trading\_day(year, month, holidays)}~;
\texttt{calendar\_month\_average}~; \texttt{storage\_floor}~;
\texttt{implied\_storage\_cost}.
\bfield{Règles} Les jours ouvrés sont les jours de semaine qui ne figurent pas dans l'ensemble des jours fériés~; un mois sans
règlement est une erreur~; en cas d'égalité dans le panier, c'est le grade classé premier par ordre alphabétique qui
l'emporte, de sorte que le résultat est déterministe.
\bfield{Tests d'acceptation} \texttt{code/firm/crude/tests/}~: l'eau à 10° API~; le
contrat de mai 2020 expirant le 21 avril 2020~; le contrat de janvier 2024 avec et sans
le jour de Noël férié~; le coût de stockage implicite du 20 avril 2020.
\bfield{Pour aller plus loin} Alimenter \texttt{firm.contracts} (livre~1, chapitre~18) avec les dates
d'expiration~; un CMA avec l'ajustement de roll pour la partie du mois valorisée sur le
second contrat.
\end{build}

\begin{omsources}
\item CFTC, \emph{Interim Staff Report on Trading in NYMEX WTI Crude Oil Futures
Contract Leading up to, on, and around April 20, 2020}, novembre 2020.
\item NYMEX Rulebook, chapter 200 (Light Sweet Crude Oil Futures), via a clearing
member's specification page~; ICE Brent Crude futures specification.
\item S\&P Global Platts, release of 8 June 2022 on WTI Midland in Dated Brent~; \emph{Specifications Guide: Europe and Africa Crude
Oil}, novembre 2025, et Brent FAQ (copies de l'Internet Archive).
\item CME Group, ``The Peaking of Storage'' (education article, 2020~; Internet Archive copy).
\item US Energy Information Administration~: glossary (API gravity)~; \emph{Today in
Energy} sur les contrats à terme Murban (2021)~; prix quotidiens des contrats à terme NYMEX, contrats 1--4.
\item IEA, \emph{Oil Market Report}, mars et août 2026.
\item United States Oil Fund, Form 8-K filings, avril 2020.
\end{omsources}

\section{Exercices}

\begin{exercise}[$\star$]\label{exo:m3:crude-oil:1}
Un brut a une densité relative de 0.87~: quelle est sa \omterm{def:m3:crude-oil:api}{densité API}~? Est-il plus léger ou plus
lourd qu'un brut à 39.6°~?
\end{exercise}

\begin{exercise}[$\star$]\label{exo:m3:crude-oil:2}
Donner le \omterm{def:m3:crude-oil:ltd}{dernier jour de négociation} du contrat WTI de décembre 2026 (le 25 novembre 2026 est
un mercredi~; ne pas tenir compte des jours fériés autres que les week-ends).
\end{exercise}

\begin{exercise}[$\star$]\label{exo:m3:crude-oil:3}
Pourquoi le \omterm{def:m3:crude-oil:brent}{Dated Brent} est-il évalué à partir d'un panier de qualités plutôt qu'à partir du seul Brent~?
\end{exercise}

\begin{exercise}[$\star\star$]\label{exo:m3:crude-oil:4}
En utilisant les barèmes d'ajustement du texte, que vaut un brut à 36° contenant 0.60~\% de soufre
face à une référence de \qty{80}{\$/\barrel} à 38° et 0.40~\%~?
\end{exercise}

\begin{exercise}[$\star\star$]\label{exo:m3:crude-oil:5}
Le stockage à Cushing coûte \qty{0.40}{\$/\barrel} par mois et le taux d'intérêt est nul. Le second
contrat se négocie à \qty{25}{\$/\barrel}. Quel est le plancher du contrat de l'échéance la plus proche~?
Que révèle un prix de 18~\$ pour l'échéance la plus proche~?
\end{exercise}

\begin{exercise}[$\star\star$]\label{exo:m3:crude-oil:6}
Le brut Arab Light pour l'Asie est facturé à la moyenne Oman/Dubai du mois de chargement, plus
le \omterm{def:m3:crude-oil:osp}{prix de vente officiel} (OSP). Avec un OSP de \qty{2.20}{\$/\barrel} et une moyenne de \qty{71.35}{\$/\barrel}, calculer
ce que paie l'acheteur pour une cargaison de 2 millions de barils~?
\end{exercise}

\begin{exercise}[$\star\star\star$]\label{exo:m3:crude-oil:7}
\emph{Programmation.} Avec \texttt{load\_wti}, compter les jours où le second contrat
dépasse de plus de 5~\$ le prix de l'échéance la plus proche, puis donner le spread le plus large
ainsi que le deuxième spread le plus large de l'historique.
\end{exercise}

\begin{exercise}[$\star\star\star$]\label{exo:m3:crude-oil:8}
\emph{Repérer l'erreur.} «~Les prix négatifs sont impossibles pour une matière première qui a un usage~:
quelqu'un prendra toujours du pétrole gratuit.~» Corriger l'affirmation.
\end{exercise}

\section{Problème~: Le dernier jour du contrat de mai}

\begin{problem}[{Problème du week-end --- 20 avril 2020 du côté acheteur}]\label{pb:m3:crude-oil:1}
Le lundi 20 avril 2020 au matin, un fonds détient 1\,000 contrats WTI de mai. Il ne dispose
d'aucun stockage à Cushing ni de capacité sur un pipeline. Le vendredi, le contrat de mai a réglé
à \qty{18.27}{\$/\barrel} et le contrat de juin à \qty{25.03}{\$/\barrel}~; les baux de stockage
ordinaires à Cushing coûtent \qty{0.40}{\$/\barrel} par mois (illustratif).

\medskip\noindent\textbf{Partie I --- La position.}
\begin{enumerate}
\item Combien de barils le fonds doit-il prendre en livraison s'il conserve sa position jusqu'à l'échéance~?
\item Quel est le \omterm{def:m3:crude-oil:ltd}{dernier jour de négociation}, et pourquoi le 20 avril est-il l'avant-dernier jour~?
\item Quelle est la valorisation au prix de marché de la position du fonds au règlement du vendredi~?
\item Quel était le spread mai--juin le vendredi, et quel coût de stockage impliquait-il~?
\item Cela était-il cohérent avec le bail habituel~?
\end{enumerate}

\medskip\noindent\textbf{Partie II --- La journée.}
\begin{enumerate}[resume]
\item Le contrat de mai est réglé à $-\qty{37.63}{\$/\barrel}$. Quelle est la
marge de variation d'une journée pour le fonds~?
\item Quel coût de stockage le spread mai--juin implique-t-il désormais~?
\item Combien de fois ce coût dépasse-t-il le bail habituel~?
\item Pourquoi le plancher de stockage n'a-t-il pas joué~?
\item Combien aurait gagné un trader disposant de 1 million de barils de capacité de stockage vide à Cushing
en achetant le contrat de mai et en vendant celui de juin au prix de règlement, en stockant pendant un mois au coût normal~?
\end{enumerate}

\medskip\noindent\textbf{Partie III --- Les alternatives.}
\begin{enumerate}[resume]
\item Combien le fonds aurait-il payé pour effectuer le roll vers juin aux prix de règlement du vendredi~?
\item Combien perd par baril un fonds à position acheteuse seule lorsqu'il effectue un roll en contango, au spread du vendredi~?
\item Pourquoi détenir des échéances plus éloignées protégeait-il un fonds~?
\item Pourquoi la méthode de règlement (transactions sur spread dans une fenêtre de deux minutes) comptait-elle~?
\item Vers quel modèle la bourse a-t-elle fait basculer ses options, et pourquoi~?
\end{enumerate}

\medskip\noindent\textbf{Partie IV --- Jugement.}
\begin{enumerate}[resume]
\item Qui se trouvait de l'autre côté~: qui peut acheter un contrat à $-\$37.63$~?
\item Que révèle cette journée sur la différence entre un indice de référence réglé en espèces et un
indice de référence à livraison physique~?
\item Un contrat à terme sur le Brent serait-il devenu négatif le même jour~? Pourquoi~?
\item Énoncer le \emph{résultat nommé}~: le coût de stockage que le marché impliquait le 20 avril.
\item En une phrase~: de quoi un prix de contrat à terme à livraison physique est-il le prix~?
\end{enumerate}
\end{problem}

\section{Questions d'entretien}

\begin{interviewq}[$\star$ \iqroles{trader}]\label{iq:m3:crude-oil:1}
Quelle est la différence entre le \omterm{def:m3:crude-oil:brent}{Dated Brent} et le prix du contrat à terme sur le Brent~?
\end{interviewq}

\begin{interviewq}[$\star$ \iqroles{trader, chercheur}]\label{iq:m3:crude-oil:2}
Comment le prix d'une matière première peut-il être négatif~?
\end{interviewq}

\begin{interviewq}[$\star\star$ \iqroles{chercheur}]\label{iq:m3:crude-oil:3}
Vous effectuez un backtest d'une stratégie sur le WTI de l'échéance la plus proche, sur une série continue, et trouvez une transaction le 20
avril 2020 qui a rapporté une fortune. Que vérifiez-vous~?
\end{interviewq}

\begin{interviewq}[$\star\star$ \iqroles{trader}]\label{iq:m3:crude-oil:4}
Le spread WTI de l'échéance la plus proche vaut $-\$1.50$ et un mois de stockage à Cushing coûte 0.40~\$. Quelle
transaction ce constat suggère-t-il, et qu'est-ce qui peut mal tourner~?
\end{interviewq}

\begin{interviewq}[$\star\star$ \iqroles{développeur, risque}]\label{iq:m3:crude-oil:5}
Votre bibliothèque de pricing calcule les prix d'options avec la formule de Black appliquée au logarithme du
prix à terme. Qu'est-ce qui cesse de fonctionner quand les contrats à terme deviennent négatifs, et que faites-vous~?
\end{interviewq}

\begin{interviewq}[$\star\star\star$ \iqroles{trader, risque}]\label{iq:m3:crude-oil:6}
Pourquoi un indice de référence fondé sur un panier de qualités et sur plusieurs marchés résiste-t-il mieux aux corners
qu'un contrat physique à qualité unique, et quel en est le coût~?
\end{interviewq}
